\documentclass[11pt,letterpaper]{article}
\usepackage[top=0.55in, bottom=0.4in, left=0.7in, right=0.7in]{geometry}
\usepackage{lmodern}
\usepackage{hyperref}
\hypersetup{colorlinks=false, urlcolor=black, pdfborder=0 0 0}
\setlength{\parindent}{0pt}
\setlength{\parskip}{0.52em}
\raggedright
\pagestyle{empty}

\begin{document}

\noindent
{\bfseries Ke-Cheng Hsu (Clio)}\\
Ph.D.\ Candidate, Department of Economics, HKUST\\
(852) 6840\,9753 \quad $|$ \quad \href{mailto:khsu@connect.ust.hk}{khsu@connect.ust.hk}\\
\href{https://cliohsu1997.github.io/}{cliohsu1997.github.io} \quad $|$ \quad \href{https://github.com/cliohsu1997}{github.com/cliohsu1997}

\vspace{0.3em}
\noindent
16 November 2026

\vspace{0.2em}
\noindent
Hiring Team\\
Asia Focus Recruitment\\
Re: Management Trainee Programme 2026 (Finance / Insurance; 1-year fast track)

\vspace{0.08em}
Dear Hiring Team,

I am writing to apply for the \textbf{Management Trainee Programme 2026}. I am a Ph.D.\ candidate in Economics at HKUST (expected June 2027; available as early as February 2027), a Hong Kong permanent resident, a native Mandarin speaker, professionally fluent in English, with basic Cantonese. I hold a bachelor's and above, with limited industry work experience (well under two years), and I want a structured 1-year path into finance with a long-term career in insurance.

The programme design matches what I am looking for. The posting describes task-oriented classroom and experiential training, job rotations and cross-functional attachment, senior mentorship, corporate-event and project-management exposure, and Mainland China / overseas business experience in a fast-paced global finance setting. I want that sequence: build financial knowledge and personal effectiveness quickly, rotate across functions to see how the business runs, take ownership of projects under a senior mentor, and grow toward leadership---with insurance as the long-term industry focus.

My background maps to the requirements. I completed a B.Sc.\ in Mathematics and Economics with a minor in \textbf{Actuarial Mathematics} at HKUST (CGA 3.9/4.3; Academic Achievement Medal, top 1\%; Dean's List), an M.Phil.\ in Economics (CGA 4.1/4.3), and I am finishing a Ph.D.\ in empirical industrial organization. The actuarial minor gave me a foundation in risk, pricing logic, and long-horizon financial products---the same language insurance businesses use. I am self-motivated and responsible: my sole-authored job market paper framed a business trade-off between sales volume and return costs, quantified options, and produced a clear recommendation pack (a moderate restocking-fee menu raises social welfare by about 2 million RMB, roughly a 5\% gain, and seller profit by about 1.2 million RMB relative to status-quo policies). With Professors Kawaguchi and Kawai I delivered weekly project updates on pricing and inventory under demand uncertainty on a fixed deadline---analytical problem-solving with ownership.

Leadership, teamwork, and communication are already part of how I work. As Teaching Assistant Coordinator and as a lecturer/TA at HKUST (4.53/5 and 4.72/5), I coordinate across people, present clearly to mixed audiences in English and Mandarin, and take responsibility for deliverables---the interpersonal side of rotations, client-facing work, and corporate projects. I bring a creative, evidence-based mindset: I turn messy data into decision-ready options rather than one-off opinions. I am confident in a fast-track trainee setting and committed to building a long-term career in the insurance industry through this programme. My resume is attached. Thank you for your consideration.

\vspace{0.3em}
Sincerely,\\[0.25em]
Ke-Cheng Hsu (Clio)

\end{document}
